\section{Future Directions\label{future}}
This section will discuss emerging frontiers for TKGQA,  aiming to stimulate further research in this field.

\subsection{More Question Types}
While existing datasets already cover some of the temporal questions, there are still more questions to be explored in the real world. 
\begin{inparaenum}[(i)]
\item  More combination of existing question types: ``Who was the \textbf{first} person to win a medal \textbf{during} the 2024 Olympic Games?'' 
\item More time granularity: Some questions demand more fine-grained granularities, such as ``When was the Long March 1 launched?'' 
\item Questions must consider the posed time: ``Where are the Seneca Indians now?''~\cite{jiaComplexTemporalQuestion2021,liskaStreamingQABenchmarkAdaptation2022a} 
\item  Predicting the future questions: ``Will the Palestinian-Israeli conflict end next year?''~\cite{jinForecastQAQuestionAnswering2021, dingForecastingQuestionAnswering2022, dingForecastTKGQuestionsBenchmarkTemporal2022} \item Common sense temporal questions: ``How often are the Olympics held?''
\end{inparaenum}

\subsection{Robust Model}
The current TKGQA methods are mostly based on the gold annotated entities and temporal signals given in the dataset~\cite{saxenaQuestionAnsweringTemporal2021,jiaComplexTemporalQuestion2021, neelamBenchmarkGeneralizableInterpretable2022}. However, these models lack robustness against real-world data, which stems from two aspects: on the one hand, the annotated entities and temporal signals in real-world data may contain errors, and the models need to have fault tolerance; on the other hand, the models need to have the ability to generalize to unseen entities and temporal signals in the training set. These limitations can be addressed in future work.

\subsection{Multi-modal TKGQA} Current TKGQA systems mainly handle plain text input. However, we experience the world with multiple modalities (e.g., video). Therefore, building a multi-modal TKGQA system that can handle multiple modalities is an important direction to investigate~\cite{yuProphetPromptingLarge2023}. 
A non-trivial challenge is effectively making a multimodal feature alignment, and complementary to better understand the temporal part.
\subsection{LLM for TKGQA} Recently, LLMs have gained significant attention for their remarkable performance across a wide range of Natural Language Processing (NLP) tasks~\cite{touvronLlamaOpenFoundation,openaiGPT4TechnicalReport2024,geminiteamGeminiFamilyHighly2024}. Existing research has also explored applying LLMs in KGQA scenarios, employing both few-shot and zero-shot learning paradigms~\cite{nieCodeStyleContextLearning2024,sunThinkGraphDeepResponsible2024, jiangStructGPTGeneralFramework2023a,baekKnowledgeAugmentedLanguageModel2023,liFewshotIncontextLearning2023a,liUnlockingTemporalQuestion2023a}
    Several emerging opportunities could further enhance the capabilities of LLMs in TKGQA systems:
    \begin{enumerate}
        \item \textbf{Multi-Agent Collaboration Interactive Reasoning for TKGQA.} 
    Recent LLM works have shifted the focus from traditional NLP tasks to exploring language agents in simulation environments that mimic real-world scenarios~\cite{zhangTimeArenaShapingEfficient2024}. Qian~\cite{qianScalingLargeLanguageModelbasedMultiAgent2024} investigates interactive reasoning and collective intelligence in autonomously solving complex
    problems. This may be further explored for temporal reasoning in temporal questions.    
    \item \textbf{Diverse Data Generation.} Numerous studies have demonstrated the effectiveness of large models in data generation~\cite{chungIncreasingDiversityMaintaining2023}, which can be used to enhance the diversity of the TKGQA dataset.
    \item \textbf{Supplementing Knowledge.} The language model itself can serve as a TKG as demonstrated by Dhingra~\cite{dhingraTimeAwareLanguageModels2022}. Additionally, LLMs possess temporal commonsense~\cite{chuTimeBenchComprehensiveEvaluation2023}, which is often absent in traditional temporal knowledge graphs. This temporal knowledge can complement existing TKGs for TKGQA.
    \end{enumerate}






